% generated by tools/bronstein_paper_table.py from tests/bronstein_paper/examples.m -- do not edit
\begin{longtable}{>{\raggedright\arraybackslash}p{0.06\textwidth}>{\raggedright\arraybackslash}p{0.40\textwidth}>{\raggedright\arraybackslash}p{0.17\textwidth}>{\raggedright\arraybackslash}p{0.22\textwidth}r}
\toprule
Id & Input & Result (\S, version) & Before & s \\ \midrule
\endhead
P0-1 & \code{EIQ}: \code{Cos[E\textasciicircum{}x]}\newline{\footnotesize the lower monomial E\textasciicircum{}x reached the base-field solver in surface form; the empty result was read as a theorem} & \code{False} (--, v0.6545) & {\footnotesize False, by accident: the base solver could not form the equation} & 0.01 \\
P0-2 & \code{EIQ}: \code{E\textasciicircum{}x*Cos[E\textasciicircum{}x]}\newline{\footnotesize the same tower, elementary: Sin[E\textasciicircum{}x]} & \code{True} (--, v0.6545) & {\footnotesize True} & 0.00 \\
P1-1 & \code{Scope}: \code{Sqrt[Log[x] + Sqrt[Log[x]]]}\newline{\footnotesize nested radicals: outside one simple radical extension, reported as such} & implementation incomplete (--, v0.6545) & {\footnotesize unevaluated, silent} & 0.00 \\
P1-2 & \code{Scope}: \code{Sqrt[x] + Sqrt[x + 1]}\newline{\footnotesize Z/2 x Z/2: a genuine compositum} & implementation incomplete (--, v0.6545) & {\footnotesize unevaluated, silent} & 0.00 \\
P2-1 & \code{RDE}: \code{\{x\textasciicircum{}(-1), 2*(Log[x]/x\textasciicircum{}2)\}}\newline{\footnotesize the solution has degree dc - da + 1 in Log[x]: the bound an off-by-one first got wrong} & \code{Log[x]\textasciicircum{}2/x} (A2, v0.6546) & {\footnotesize undecided} & 0.00 \\
E6b-1 & \code{(Sqrt[x] + (x + 1)\textasciicircum{}(1/3))/(x*(x + 1))}\newline{\footnotesize y\textasciicircum{}6 = x\textasciicircum{}3 (x+1)\textasciicircum{}2, y = Sqrt[x] (x+1)\textasciicircum{}(1/3); answer in the user's radicals} & elementary (E6b, v0.6547) & {\footnotesize implementation incomplete} & 0.22 \\
E6b-2 & \code{Sqrt[x]*(x + 1)\textasciicircum{}(1/3)}\newline{\footnotesize Chebyshev's binomial case, certified by the Risch-Trager-Bronstein machinery on y\textasciicircum{}6 = x\textasciicircum{}3 (x+1)\textasciicircum{}2} & not elementary (certified) (E6b, v0.6547) & {\footnotesize implementation incomplete} & 0.02 \\
E6b-3 & \code{D[Sqrt[x - 1]*(x + 1)\textasciicircum{}(1/3)*E\textasciicircum{}x, x]}\newline{\footnotesize the compositum at the bottom of a transcendental tower (Shape B, E\textasciicircum{}x on top)} & elementary (E6b, v0.6547) & {\footnotesize implementation incomplete} & 0.17 \\
E6b-5 & \code{Sqrt[x]*(x + 1)\textasciicircum{}(1/3)*Log[x]}\newline{\footnotesize Log[x] over the compositum: certified by the coefficient Risch DE over the base} & not elementary (certified) (E6b, v0.6547) & {\footnotesize implementation incomplete} & 0.02 \\
E6b-6 & \code{(x*Sqrt[Log[x]]*(Log[x] + 1)\textasciicircum{}(1/3))\textasciicircum{}(-1)}\newline{\footnotesize Chebyshev's case in u = Log[x] (u\textasciicircum{}(-1/2) (1 + u)\textasciicircum{}(-1/3)), decided in the tower} & not elementary (certified) (E6b, v0.6547) & {\footnotesize implementation incomplete} & 0.03 \\
E6b-4 & \code{Sqrt[Log[x]]*((Log[x] + 1)\textasciicircum{}(1/3)/x)}\newline{\footnotesize the compositum over the transcendental tower Q(x)(Log[x])} & not elementary (certified) (E6b, v0.6547) & {\footnotesize implementation incomplete} & 0.03 \\
E6d-1 & \code{Sqrt[x\textasciicircum{}2*(x + 1)]/x\textasciicircum{}2 + (Sqrt[x + 1])\textasciicircum{}(-1)}\newline{\footnotesize Sqrt[x\textasciicircum{}2 (x+1)] = eps x Sqrt[x+1], eps = +-1 locally constant; the answer is a sum over sign regions} & elementary (E6d, v0.6549) & {\footnotesize implementation incomplete} & 0.08 \\
E6d-2 & \code{Sqrt[x\textasciicircum{}2*(x\textasciicircum{}3 + 1)] - x*Sqrt[x\textasciicircum{}3 + 1]}\newline{\footnotesize identically 0 where eps = 1, the elliptic -2 x Sqrt[x\textasciicircum{}3+1] where eps = -1} & not elementary (certified) (E6d, v0.6549) & {\footnotesize implementation incomplete} & 0.01 \\
E6d-3 & \code{Sqrt[x\textasciicircum{}2*(x\textasciicircum{}3 + 1)]*(Sqrt[x\textasciicircum{}3 + 1]/x)}\newline{\footnotesize the product is eps (x\textasciicircum{}3 + 1): rational on each region} & elementary (E6d, v0.6549) & {\footnotesize implementation incomplete} & 0.04 \\
E6d-4 & \code{Sqrt[Log[x]\textasciicircum{}2*(Log[x] + 1)] + Sqrt[Log[x] + 1]/x}\newline{\footnotesize over the tower: Log[x] Sqrt[Log[x] + 1] on the region eps = 1} & not elementary (certified) (E6d, v0.6549) & {\footnotesize implementation incomplete} & 0.02 \\
E6d-5 & \code{(1 + E\textasciicircum{}x\textasciicircum{}2)/2 + ((1 - E\textasciicircum{}x\textasciicircum{}2)/2)*Sqrt[x + 1]*Sqrt[(x + 1)\textasciicircum{}(-1)]}\newline{\footnotesize equals 1 on every open set; the sign -1 occurs only on the cut x < -1 and must never certify} & implementation incomplete (E6d, v0.6550) & {\footnotesize not elementary (v0.6549: false, a sign seen only on a branch cut)} & 0.01 \\
E6c-1 & \code{E\textasciicircum{}(x + Log[x]/2)/x}\newline{\footnotesize E\textasciicircum{}x/Sqrt[x] in disguise} & not elementary (certified) (E6c, v0.6547) & {\footnotesize undecided (a new hyperexponential monomial whose Risch DE needs a radical)} & 0.00 \\
E6c-4 & \code{E\textasciicircum{}(x\textasciicircum{}2 + Log[Log[x]]/2)}\newline{\footnotesize E\textasciicircum{}(x\textasciicircum{}2) Sqrt[Log[x]]: Shape A with the exponential below the radicand's monomial} & not elementary (certified) (E6c, v0.6547) & {\footnotesize undecided} & 0.01 \\
E6c-2 & \code{D[E\textasciicircum{}(x\textasciicircum{}2 + Log[x + 1]/3), x]}\newline{\footnotesize E\textasciicircum{}(x\textasciicircum{}2) (x+1)\textasciicircum{}(1/3): a radical at the bottom under an exponential} & elementary (E6c, v0.6547) & {\footnotesize undecided} & 0.01 \\
E6c-3 & \code{D[E\textasciicircum{}(Log[Log[x]]/2 + x), x]}\newline{\footnotesize Log[Log[x]] is rewritten away and must be pruned} & elementary (E6c, v0.6547) & {\footnotesize implementation incomplete (radical in the middle of the tower)} & 0.01 \\
E6a-1 & \code{Log[x\textasciicircum{}2]*(Sqrt[Log[x]]/x)}\newline{\footnotesize answer carries the branch difference Log[x\textasciicircum{}2] - 2 Log[x] explicitly} & elementary (E6a, v0.6547) & {\footnotesize implementation incomplete} & 0.01 \\
E6a-2 & \code{(Log[x\textasciicircum{}2] - 2*Log[x])*Sqrt[Log[x]]}\newline{\footnotesize zero on Re x > 0; certified by Hermite-Lindemann + specialisation (Thm E6a(3))} & not elementary (certified) (E6a, v0.6547) & {\footnotesize implementation incomplete} & 0.00 \\
E6a-3 & \code{Log[x*(x + 1)]*Sqrt[Log[x] + Log[x + 1]]*(x\textasciicircum{}(-1) + (x + 1)\textasciicircum{}(-1))}\newline{\footnotesize a product logarithm over two log monomials} & elementary (E6a, v0.6547) & {\footnotesize implementation incomplete} & 0.07 \\
E6a-4 & \code{(Log[E\textasciicircum{}x] - x)*Sqrt[Log[x]] + x\textasciicircum{}(-1)}\newline{\footnotesize the branch difference vanishes only on |Im x| < Pi; certified at a sample with large imaginary part} & not elementary (certified) (E6a, v0.6547) & {\footnotesize Log[x] from the Automatic cascade (wrong: PossibleZeroQ[Log[E\textasciicircum{}x] - x] is True)} & 0.01 \\
E3-1 & \code{(x*(1 + Log[x]\textasciicircum{}3)\textasciicircum{}(1/3))\textasciicircum{}(-1)}\newline{\footnotesize genus 1 (Fermat cubic over Q(x)(Log[x])); residues at theta = infinity} & elementary (E3, v0.6548) & {\footnotesize implementation incomplete} & 2.41 \\
E3-2 & \code{E\textasciicircum{}x/(1 + E\textasciicircum{}(3*x))\textasciicircum{}(1/3)}\newline{\footnotesize the exponential analogue: special places over theta = 0 and infinity} & elementary (E3, v0.6548) & {\footnotesize implementation incomplete} & 2.40 \\
E3-3 & \code{(x*(1 + Log[x]\textasciicircum{}4)\textasciicircum{}(1/4))\textasciicircum{}(-1)}\newline{\footnotesize m = 4, genus 3} & elementary (E3, v0.6548) & {\footnotesize implementation incomplete} & 5.68 \\
E3-4 & \code{D[(1 + Log[x] + Log[x]\textasciicircum{}3)\textasciicircum{}(1/3) + Log[Log[x] + (1 + Log[x] + Log[x]\textasciicircum{}3)\textasciicircum{}(1/3)], x]}\newline{\footnotesize algebraic part plus a log part with finite and special residues on y\textasciicircum{}3 = L\textasciicircum{}3 + L + 1} & elementary (E3, v0.6548) & {\footnotesize implementation incomplete} & 1.80 \\
E3-5 & \code{(1 + Log[x])/(x*(1 + Log[x] + Log[x]\textasciicircum{}3)\textasciicircum{}(1/3))}\newline{\footnotesize the residue divisor has infinite order (orders mod 13 and 37 disagree)} & not elementary (certified) (E3, v0.6548) & {\footnotesize implementation incomplete} & 1.18 \\
E3-6 & \code{D[x + Log[-E\textasciicircum{}x + (1 + E\textasciicircum{}(3*x))\textasciicircum{}(1/3)], x]}\newline{\footnotesize exponential tower, genus 1} & elementary (E3, v0.6548) & {\footnotesize implementation incomplete} & 4.93 \\
E41-1 & \code{Sqrt[x + I]/x}\newline{\footnotesize a curve over Q(i)(x); FriCAS agrees} & elementary (E4.1, v0.6551) & {\footnotesize implementation incomplete} & 0.07 \\
E41-2 & \code{D[Log[E\textasciicircum{}(I*x) + Sqrt[1 + E\textasciicircum{}(2*I*x)]], x]}\newline{\footnotesize Shape A over the exponential E\textasciicircum{}(I x): an inverse-hyperbolic of E\textasciicircum{}(I x)} & elementary (E4.1, v0.6551) & {\footnotesize implementation incomplete} & 0.09 \\
E41-3 & \code{(Sqrt[E\textasciicircum{}(2*I*x) + 1])\textasciicircum{}(-1)}\newline{\footnotesize residues at the special places of E\textasciicircum{}(I x) over Q(i)} & elementary (E4.1, v0.6551) & {\footnotesize implementation incomplete} & 0.09 \\
E41-4 & \code{(1 + I)*(Sqrt[x]/(x\textasciicircum{}2 + I))}\newline{\footnotesize residue field Q(i)(zeta\_8): i adjoined and written in the primitive element} & elementary (E4.1, v0.6551) & {\footnotesize implementation incomplete} & 0.09 \\
E41-5 & \code{Sqrt[Log[x] + I]}\newline{\footnotesize FriCAS: implementation incomplete (constant residues); Risch DE b' + b (1 + 1/(2 (u + I))) = 1 has no rational solution} & not elementary (certified) (E4.1, v0.6551) & {\footnotesize implementation incomplete} & 0.01 \\
\bottomrule
\end{longtable}
